% !TeX root = ../papers/nc-p3.tex
\section{二次元単純を印した有限窓の全階数型}
\label{sec:two-dimensional-marked-tensor-forms}
\begin{theorem}\label{thm:two-dimensional-marked-tensor-forms}
四辺が有限な標準有限Frobenius窓に、共通次数倍率を除いた頂点次元二の単純を一つ印す。
各文字作用 $A_i:V_i\to\operatorname{Hom}(U_i,U_{i+1})$ の階数3の位置集合を
$D\subset\mathbb Z/4$ とする。この印付き窓の候補を次の線型空間で表示できる。
$D$ が全位置である場合だけ追加の共役類 $h\in\mathrm{PGL}_2$ を用いる。
空間の次元は通常81で、全位置が階数3かつ $h=1$ の場合だけ88である。
その中で隣接二切断が階数6となることと標準有限Frobenius・有限辺の検査を課す。
この表示上で指定した文字・頂点の基底変更で移り合うことが印付き添字保持同型の必要十分条件である。
これには標準AS性をまだ判定していない窓も含まれる。
\end{theorem}
\begin{proof}
階数の二択と、階数3なら像の零化線が可逆な行列で張られることは
定理\ref{thm:finite-two-dimensional-quadrics}による。
この線は $U_i$ と $U_{i+1}$ の射影同定を定める。階数4の位置では文字空間を
$U_i^*\otimes U_{i+1}$、作用を恒等とする。階数3の位置では
一文字 $e_i$ を核とし、残りを上の同定による跡零の三空間とし、作用はその上への射影とする。
不完全な辺の連鎖が閉じないときそれらの同定は頂点基底で全て恒等になる。
全て階数3なら三同定を恒等とし、最後に一周の合成 $h$ の共役類だけが残る。

四テンソル $w$ を、この分解のもと $e_i$ を選んだ位置がちょうど $I\subset D$ である
各成分 $w_I$ に分ける。各残存位置の二次元因子二つを、順に入口と出口と呼ぶ。
関係 $R_i$ は $w$ の他の二位置での全縮約である。
従って二辺の積の作用が零となる条件は、両位置が残るとき
出口と次の入口の縮約が零という線型条件である。殺された位置が隣にある条件は消える。
これらの条件は各 $I$ で独立である。残った階数3位置では、入口・出口同士の縮約も
跡零条件で零になる。二次元の交代体積で双対を同定すると、一つの縮約が零であることは
その二つの因子の対称性そのものである（一次元の行列式因子は基底で自明化する）。

従って二因子をこの対称条件ごとに結び、連結する鎖に $m$ 個の因子があれば、
その部分は $\operatorname{Sym}^m(k^2)$ になる。隣接互換が全対称群を生成するからである。
鎖ごとに二次元空間を同定してよく、その次元の積 $\prod(m+1)$ がこの $I$ の寄与である。
唯一の閉路は $D$ が全体で $I=\varnothing$ の場合である。このとき最後の辺以外で
対称八乗を得た後、残る条件は二因子への作用 $h\otimes h^{-1}$ で不変なことになる。
実際最後の互換を片方だけの $h$ で共役し、既に通常の互換で不変であることを使えばよい。
$h=1$ なら次元9である。二固有値なら両固有文字が混じる係数が消えて純八乗二つが残る。
非自明な単冪なら $h=1+N$, $Ny=x$ の基底を選ぶと条件は $N_1w=N_2w$ である。
両因子が $y$ の成分だけがこの差に寄与するので、$x^8,x^7y$ の二次元が残る。
全16個の $D$ とそれらの各 $I$ の和を取ると上の81,88となる。
\path{computations/verify_two_dimensional_marked_forms.py} は鎖の数え上げと八乗上の二条件を検算する。

この空間上で得たテンソルを元の二切断の階数検査へ戻せば、関係は必ず指定した
$A_{i+1}A_i$ に消される。各 $A_i$ の像はどの非零ベクトル上でも全射なので、
得られる折り畳み表現は単純である。
残る同値は、四頂点基底 $g_i$ と四文字基底 $f_i$ で
$A_i f_i=g_{i+1}A_i g_i^{-1}$ を満たし、$\bigotimes f_i$ で $w$ の直線を運ぶものである。
頂点基底は階数3の辺で射影同定を運び、各 $f_i$ はそこで核文字の任意の非零倍率と
三空間から核への任意のせん断を許す。共通次数倍率も許すが、全ての $f_i$ の
共通スカラーを変えることで上の条件に吸収できる。各辺の倍率を変えた表示も、
四頂点倍率と共通次数倍率に分解できる。これで規格化してからの余分な同一視も全て記述した。
\end{proof}

\begin{corollary}\label{cor:two-dimensional-full-arrow-apolarity}
全ての作用が階数4なら、表示は四つの二元二次形式のテンソル
$F\in\bigotimes_i\operatorname{Sym}^2 U_i$ の射影化で、残る群は
$\prod_i\mathrm{PGL}(U_i)$ である。二つの階数6条件は多重次数
$(1,2,1,0),(0,1,2,1)$ の十二行十二列の偏極行列がともに階数6であることに一致する。
この階数で有限窓を与える $F$ は、多重二次のSegre--Veronese像の長さ6以下の
零次元部分スキームの線型包には入らない。
\end{corollary}
\begin{proof}
四つの鎖は各頂点の出口・入口の一組ずつなので二乗を取る。最初の二文字をまとめると、
両端の二次元因子を一つずつ、中間の頂点は対称二乗全部を取る。反対側も同様である。
二つの端で対称二乗を一因子ずつに分ける偏極が元の行列切断を与え、
余分の交代成分は行・列とも零である。これで表示サイズと多重次数が従う。
もし長さ6以下の $Z\subset(\mathbb P^1)^4$ の像が $F$ を張れば、
この切断は $Z$ への左右の評価とその上の一つの線型汎関数による積対合を通る。
階数6より $Z$ は長さ6で両評価も全射、行の像は評価双対の像である。
$Z$ のどの台点でもその評価は左二文字の可分解テンソルであり、関係六空間に入る。
これは有限Frobenius辺で純関係がないこと（定理\ref{thm:finite-cubic-separators}の証明）に反する。
\end{proof}
